\documentclass[float=false, crop=false]{standalone}
\usepackage{hyperref}
\usepackage{comment}

%\oddsidemargin 1.0in
\oddsidemargin 0.0in
\topmargin -0.1in

\pagestyle{empty}
\textheight 8.5in
\textwidth 6.5in

\begin{document}
%-----------------------------------------------------
\vspace*{-0.9in}
\begin{huge}
\centerline{\bf Indian Institute of Science}
\end{huge}
%-----------------------------------------------------
\vskip 0.08in
\begin{large}
\noindent
{\sf Department of Computational and Data Sciences}   \hfill
{\sf Office: +91-80-2293 2496 }\\
%
{\sf  Bangalore, India - 560012}   \hfill
{\sf  FAX: +91-80-2360 2648}  \\
\hrule
\end{large}
\begin{flushright}
\today \\
\end{flushright}

\begin{large}
%\vskip 0.1in
\noindent
To\\
The Editor\\
Medical Physics \\

%\hfill


%------------------------------>>>>>>   INPUT DATE

%\vskip 0.1in
\noindent
Dear Editor,

%\newcommand{\pp}{\vskip 0.10in}
\par
We would like to submit a majorly revised manuscript titled `\emph{Performance Benchmarking of Deep Learning Models for Real-Time Median Nerve Segmentation and Cross-Sectional Area Measurement in Ultrasound Imaging}' to be considered for Medical Physics. \\


\noindent
The main contributions of this work:

\noindent
1. Proposed an efficient, lightweight CNN model called MNSeg-Net for fully automated segmentation of the median nerve from a single ultrasound frame, supporting real-time inference at 43 frames per second.

\noindent
2. The model can efficiently segment the median nerve from wrist to elbow and calculate cross-sectional area (CSA), matching the performance of expert sonographers.

\noindent
3. Clinically validated the proposed model by testing on clinical data from normal subjects and patients with mild and severe carpal tunnel syndrome.

\noindent
4. Developed an end-to-end clinical system with a Python-based graphical user interface for real-time median nerve segmentation to assist sonographers. This setup is already in use for research at Aster-CMI Hospital.

\noindent
5. Introduced a novel lightweight architecture combining a redesigned U2Net as the main network with a multi-scale feature fusion sub-network module to effectively capture contextual information and enhance feature representation.

\noindent
6. Achieved state-of-the-art segmentation performance (94.7\% DSC at wrist, 83.4\% DSC from wrist to elbow) with only 2.46 million parameters, matching the accuracy of much larger models while being significantly more efficient.

\noindent
7. Demonstrated the model's ability to process up to 32 frames per second in a clinical setting with the developed system, including resizing, segmentation, contour generation, and CSA calculation. The developed codes are available as an open source for enthusiastic users at \href{https://github.com/venkateshvaddadi/Real-Time-Median-Nerve-Segmentation/}{https://github.com/venkateshvaddadi/Real-Time-Median-Nerve-Segmentation/}\\

This work has not been submitted (or is under consideration) to any other publication. Please do not hesitate to contact me if you require further information on this submission.\\


\noindent Sincerely,\\
Phaneendra K. Yalavarthy, Ph.D. (Corresponding Author)\\
Computational and Data Sciences\\ 
Indian Institute of Science, Bangalore - 560012 India. \\
Tel. No: +91-80-2293 2496; FAX: +91-80-2360 2648.\\
E-mail: yalavarthy@iisc.ac.in\\
Web: https://cds.iisc.ac.in/faculty/yalavarthy/
%\hspace*{3.0in}   \\
\end{large}
\end{document}
